% !TeX program = xelatex
\documentclass[11pt,a4paper]{article}

\usepackage{kotex}
\usepackage{geometry}
\geometry{left=26mm,right=26mm,top=27mm,bottom=28mm}
\usepackage{graphicx}
\usepackage{booktabs}
\usepackage{longtable}
\usepackage{tabularx}
\usepackage{array}
\usepackage{multirow}
\usepackage{caption}
\usepackage{subcaption}
\usepackage[table]{xcolor}
\usepackage{enumitem}
\usepackage{fancyhdr}
\usepackage{hyperref}
\usepackage{setspace}
\usepackage{float}
\usepackage{amsmath}
\usepackage{amssymb}
\usepackage{threeparttable}
\usepackage{pdflscape}

\hypersetup{
    colorlinks=true,
    linkcolor=black,
    urlcolor=black,
    citecolor=black,
    pdfauthor={고려대학교 SDA},
    pdftitle={SSG 랜더스 외국인 선수 영입 전략},
    pdfsubject={Sports Data Mining Report},
    pdfkeywords={KBO, SSG Landers, foreign player scouting, data mining, ensemble model}
}

\newcolumntype{L}[1]{>{\raggedright\arraybackslash}p{#1}}
\newcolumntype{C}[1]{>{\centering\arraybackslash}p{#1}}
\newcommand{\tight}{\setlength{\itemsep}{0.25em}\setlength{\parskip}{0pt}}
\newcommand{\code}[1]{\texttt{#1}}
\newcommand{\figpath}[1]{assets/#1}

\captionsetup{font=small,labelfont=bf}
\setstretch{1.18}
\setlist[itemize]{leftmargin=1.5em}
\setlist[enumerate]{leftmargin=1.7em}

\pagestyle{fancy}
\fancyhf{}
\lhead{\small SSG 랜더스 외국인 선수 영입 전략}
\rhead{\small \today}
\cfoot{\thepage}

\begin{document}

\begin{titlepage}
    \centering
    \vspace*{42mm}
    {\Huge\bfseries SSG 랜더스 외국인 선수 영입 전략\par}
    \vspace{9mm}
    {\Large 데이터마이닝 기반 외국인 타자 및 선발투수 후보 선정 보고서\par}
    \vspace{14mm}
    {\large Data: KBO/STATIZ 2023-2026 \quad | \quad MLB/MiLB 2023-2026 \quad | \quad Candidate pool Hitter 736 / Pitcher 1,009\par}
    \vfill
    {\large 고려대학교 SDA\par}
    \vspace{6mm}
    {\large June 23, 2026\par}
\end{titlepage}

\pagenumbering{roman}

\section*{초록 (Abstract)}
\addcontentsline{toc}{section}{초록 (Abstract)}

\subsection*{국문}
본 보고서는 SSG 랜더스의 외국인 타자 및 외국인 선발투수 영입 문제를 하나의 데이터마이닝 의사결정 문제로 재구성한다. 목표는 성적이 좋은 선수를 단순 나열하는 것이 아니라, SSG의 반복 약점을 정량적으로 정의하고, 그 약점을 완화하면서 KBO 번역 가능성과 시장 접근성을 동시에 갖춘 후보를 찾는 것이다. 최종 검증 단계에서는 모델 지지 신호, 시장 접근성, 계약/비용 조건, 의료 확인 필요성을 함께 반영하는 통합 모델로 재구성했다. KBO/STATIZ, MLB Savant, MiLB 성적, roster/transaction, public salary/contract signal, medical status, 기존 모델 산출 테이블을 통합했고, Team Need Mining, Historical KBO Success/Failure, Similarity/Archetype, KBO Translation, 시장 접근성 모형, 의료/실행 가능성 검증 모형을 앙상블했다. 검증 결과 타자 최종 후보 3인은 Will Brennan, Dominic Fletcher, Dylan Carlson으로 수정했고, Luis Matos는 초기 발견 모델에서 강한 후보였지만 age 24와 MLB 재도전 가치 때문에 접촉 보류로 재분류했다. 투수 최종 후보 3인은 Josh Fleming, Keegan Thompson, Kolby Allard이며, 최종 접촉 1순위는 Will Brennan과 Josh Fleming이다.

\subsection*{English}
This report reframes SSG Landers' foreign hitter and starting pitcher acquisition as a structured data-mining decision problem. Rather than ranking surface statistics, the analysis identifies SSG's repeated game-state weaknesses, converts them into player-level feature requirements, and selects candidates with both KBO translation potential and practical market access. In the final validation layer, the decision rule combines model support, market access, contract/cost constraints, and medical checks before contact priority is assigned. The revised hitter board is Will Brennan, Dominic Fletcher, and Dylan Carlson, while Luis Matos is moved to an early-screen hold because his age and remaining MLB reset value create player-side and club-rights uncertainty. The pitcher board is Josh Fleming, Keegan Thompson, and Kolby Allard. The recommended first contacts are Will Brennan and Josh Fleming.

\begin{center}
\begin{minipage}{0.90\linewidth}
\rowcolors{1}{gray!10}{gray!10}
\begin{tabularx}{\linewidth}{L{0.96\linewidth}}
\rowcolor{gray!35}
\textbf{Key Takeaways} \\
\begin{itemize}\tight
    \item SSG의 보강 문제는 단순 장타 부족이 아니라 RHP game-script, run-kill avoidance, extra-out resilience가 결합된 game-state interaction 문제다.
    \item 후보 선별은 전체 시장에서 시작해 후보 생성 모듈, 데이터마이닝 1차 검증, 현실성 검증 단계, 최종 접촉 순위로 좁히는 funnel로 설계했다.
    \item 타자는 현실성 제외 조건을 강화해 Will Brennan, Dominic Fletcher, Dylan Carlson을 최종 접촉 후보군으로 제시한다.
    \item Matos는 초기 발견 모델에서 강한 후보였지만 age 24와 MLB 재도전 가치 때문에 최종 후보가 아니라 접촉 보류가 맞다.
    \item Suwinski는 BB\%와 barrel upside가 있지만 실패 경고 신호와 K\% 위험이 커 최종 Top 3가 아니라 watchlist가 맞다.
    \item 투수는 supervised classifier만으로 확정 추천하기 어렵기 때문에 command, starter floor, 의료 확인, 시장 접근성 검증을 결합했다.
    \item 최종 접촉 1순위는 타자 Will Brennan, 투수 Josh Fleming이다.
\end{itemize}
\end{tabularx}
\end{minipage}
\end{center}

\clearpage
\tableofcontents
\clearpage
\listoftables
\clearpage
\listoffigures
\clearpage

\pagenumbering{arabic}

\section{서론}

\subsection{문제 정의}
이 보고서의 대상은 ``좋은 선수''가 아니라 ``SSG가 실제로 데려와야 하는 선수''다. 외국인 선수 영입은 순수 선수 평가와 다르다. 특정 팀의 약점, KBO 번역 가능성, 계약 가능성, 의료 상태, 비용 조건이 동시에 충족되어야 한다. 따라서 본 분석은 후보를 점수화하는 단일 랭킹 문제가 아니라, 여러 조건을 통과하는 후보를 단계적으로 좁히는 데이터마이닝 funnel 문제로 설계했다.

\subsection{핵심 지표}
야구 지표에 익숙하지 않은 독자도 결과를 해석할 수 있도록 핵심 지표를 표 \ref{tab:metric_guide}에 정리한다.

\begin{table}[H]
\centering
\caption{주요 야구 지표 요약}
\label{tab:metric_guide}
\small
\begin{tabularx}{\linewidth}{L{0.18\linewidth}L{0.32\linewidth}L{0.40\linewidth}}
\toprule
지표 & 의미 & 실무 해석 \\
\midrule
OBP & 출루율 & 이닝을 이어가는 능력, 상위타순 안정성 \\
SLG/ISO & 장타력 및 순장타력 & 출루를 득점으로 바꾸는 능력 \\
K\% / Whiff\% & 삼진 및 헛스윙 위험 & KBO 적응 실패와 run-kill 위험 \\
BB\% & 볼넷 생산성 & 공을 고르고 이닝을 늘리는 능력 \\
xwOBA & 타구 질 기반 예상 공격력 & 표면 성적보다 지속 가능성 평가에 유리 \\
BB/9 & 9이닝당 볼넷 & 투수의 traffic 생성 위험 \\
HR/9 & 9이닝당 피홈런 & 장타 허용 및 대량실점 위험 \\
GS / IP & 선발 등판 및 이닝 & 선발 지속성과 5이닝 floor 평가 \\
\bottomrule
\end{tabularx}
\end{table}

\subsection{분석 질문}
본 보고서는 표 \ref{tab:rq}의 여섯 질문을 따라 전개된다.

\begin{table}[H]
\centering
\caption{분석 질문}
\label{tab:rq}
\small
\begin{tabularx}{\linewidth}{C{0.10\linewidth}L{0.48\linewidth}L{0.31\linewidth}}
\toprule
RQ & 질문 & 최종 산출물 \\
\midrule
RQ1 & SSG의 핵심 보강 문제는 무엇인가? & SSG weakness rule \\
RQ2 & 약점은 어떤 선수 조건으로 변환되는가? & 타자/투수 feature contract \\
RQ3 & 후보군은 어떻게 생성되고 줄어드는가? & 전체 후보군에서 실행 보정 board \\
RQ4 & 각 모델은 어떤 관점에서 후보를 평가하는가? & 모델 구성표와 성능 평가표 \\
RQ5 & 타자 및 투수 최종 후보 3명은 누구인가? & slot별 실행 보정 Top 3 \\
RQ6 & 최종 접촉 1순위는 누구인가? & 접촉 우선순위 Plan A \\
\bottomrule
\end{tabularx}
\end{table}

\subsection{의사결정 구조}
최종 결정은 data-mining ranking과 현실성 반영 접촉 순위를 분리한다. Data-mining ranking은 구조화 지표가 과거 KBO 성공/실패 패턴과 얼마나 닮았는지를 뜻하고, 최종 접촉 순위는 계약, 비용, 의료, 시장 접근성, KBO행 가능성을 함께 적용한 실행 순위다. 이 구조 때문에 초기 발견 후보가 최종 접촉 1순위가 아닐 수 있다. 본 보고서에서 Luis Matos는 타자 발견 단계에서 강한 후보지만 age 24와 MLB 재도전 가치 때문에 최종 접촉 후보군이 아니라 hold로 내려가고, Will Brennan이 현실형 접촉 1순위로 올라간다.

\section{데이터 및 전처리}

\subsection{데이터 소스}
분석에 사용한 데이터의 역할은 표 \ref{tab:data_sources}와 같다. 모델 입력은 정량 structured data로 한정했다.

\begin{longtable}{L{0.23\linewidth}L{0.15\linewidth}L{0.18\linewidth}L{0.29\linewidth}}
\caption{데이터 구성 요약}
\label{tab:data_sources}\\
\toprule
데이터 소스 & 기간 & 단위 & 사용 목적 \\
\midrule
\endfirsthead
\toprule
데이터 소스 & 기간 & 단위 & 사용 목적 \\
\midrule
\endhead
KBO/STATIZ 팀/선수 데이터 & 2023-2026 & 팀, 선수, 경기, 상황 & SSG 문제 정의와 KBO label 구성 \\
SSG 상황별 경기 데이터 & 2026 & game-state rule & Team Need Mining Model \\
MLB Savant & 2023-2026 & pitch/player-season & 타자 기초 역량과 KBO translation \\
MiLB 성적 데이터 & 2025-2026 중심 & player-season/role & 후보 시장 구축과 투수 starter floor \\
MLB roster/transaction & 2025-10 이후 & 선수 상태 & 시장 접근성 모형 \\
NPB/CPBL 공개 roster/stat seed & 2026 & 선수/리그 & 시장 depth와 아시아쿼터 확장성 확인 \\
public salary/contract/medical & 2026-06 확인 & 선수 & 시장/의료 실행 검증 \\
기존 output tables & 프로젝트 산출물 & 모델/후보/검증 단계 & 최종 보고서 재현성과 후보 board 구성 \\
\bottomrule
\end{longtable}

\subsection{분석 단위}
타자는 선수-시즌 및 최근 plate appearance 단위를 기본으로 하며, 외야 및 지명타자 role, 좌우타, contact floor, on-base/damage profile을 함께 본다. 투수는 선수-시즌 및 최근 IP/GS 단위를 기본으로 하며, 선발 지속성, command stability, traffic damage control, KBO/ABS translation을 함께 본다.

\subsection{후보군 생성}
후보 시장은 MLB 40-man 경계선, DFA/outright/minor league contract 선수, AAA regular role 선수, 최근 트랜잭션 후보, KBO행 가능성이 있는 나이와 계약 상태의 선수로 구성했다. 후보 funnel은 표 \ref{tab:candidate_funnel}과 같다.

\begin{table}[H]
\centering
\caption{후보 funnel}
\label{tab:candidate_funnel}
\small
\begin{tabularx}{\linewidth}{L{0.36\linewidth}C{0.14\linewidth}C{0.14\linewidth}L{0.25\linewidth}}
\toprule
단계 & 타자 잔존 & 투수 잔존 & 판정 의미 \\
\midrule
Step 1. 전체 구조화 시장 & 736 & 1,009 & 모델 입력 가능한 후보 pool 구성 \\
Step 2. 후보 생성 모듈 & 16 & 18 & Market, Team Fit, Upside Screen 통과 \\
Step 3. 데이터마이닝 1차 검증 & 6 & 3 & 접근성, 표본, 부상 flag 1차 통과 \\
Step 4. 초기 Top 3 & 3 & 3 & 순수 모델 점수와 반복 등장 신호 \\
Step 5. 실행 보정 Top 3 & 3 & 3 & 계약, 의료, 실행 가능성 반영 \\
\bottomrule
\end{tabularx}
\end{table}

\subsection{Feature Contract}
Team Need Mining의 핵심은 팀 약점을 선수 조건으로 번역하는 것이다. 표 \ref{tab:feature_contract}는 SSG 약점과 선수 feature의 연결이다.

\begin{figure}[H]
\centering
\includegraphics[width=0.90\linewidth]{\figpath{v2_hidden_weakness_rules.png}}
\caption{SSG weakness mining rule}
\label{fig:weakness}
\end{figure}

\begin{longtable}{L{0.12\linewidth}L{0.25\linewidth}L{0.26\linewidth}L{0.25\linewidth}}
\caption{SSG 약점과 선수 feature contract}
\label{tab:feature_contract}\\
\toprule
슬롯 & SSG weakness & Player feature & 해석 \\
\midrule
\endfirsthead
\toprule
슬롯 & SSG weakness & Player feature & 해석 \\
\midrule
\endhead
타자 & RHP game-script unlocker & vs RHP on-base damage & 우투 선발 경기에서 외야/DH가 이닝을 다시 여는 능력 \\
타자 & run-kill avoidance & low GDP and low chase risk & 병살, 삼진, 약한 타구로 이닝을 끊지 않는 능력 \\
타자 & two-strike survival & two-strike contact floor & 불리한 카운트에서 최소 생산성을 유지하는 능력 \\
타자 & role stability & corner OF or DH continuity & 외야/DH 슬롯에서 바로 쓸 수 있는 역할 안정성 \\
투수 & traffic-command stabilizer & low free-pass volatility & 추가 출루 이후 볼넷으로 이닝을 키우지 않는 능력 \\
투수 & starter length floor & five-inning floor & 5이닝 이상 버틸 수 있는 최소 근거 \\
투수 & extra-out resilience & damage control after traffic & 수비 실수나 추가 출루 이후 대량실점을 막는 능력 \\
투수 & ABS adaptation & zone command & ABS 환경에서 볼넷 리스크가 커지지 않는 능력 \\
\bottomrule
\end{longtable}

\section{방법}

\subsection{전체 모델 구조}
본 보고서는 하나의 최종 앙상블 모델을 사용한다. 다만 내부 base learner는 서로 다른 질문을 담당한다. 표 \ref{tab:model_family}는 각 모델의 역할이다.

\begin{longtable}{L{0.24\linewidth}L{0.25\linewidth}L{0.23\linewidth}L{0.16\linewidth}}
\caption{앙상블 base learner 구성}
\label{tab:model_family}\\
\toprule
모델명 & 목적 & 출력 & 사용 강도 \\
\midrule
\endfirsthead
\toprule
모델명 & 목적 & 출력 & 사용 강도 \\
\midrule
\endhead
Model A. Team Need Mining & SSG 약점을 game-state interaction으로 정의 & weakness rule, feature contract & 공통 기준 \\
Model B. Historical KBO Success/Failure & 입단 전 정보로 성공/실패 패턴 유사도 추정 & 모델 지지/실패 경고 점수 & 타자 후보 비교용, 투수 진단 \\
Model C. Similarity/Archetype Matching & KBO 성공 유형 또는 SSG 필요 유형과 유사도 계산 & archetype score & 보조 \\
Model D. KBO Translation Risk & MLB/MiLB 성과가 KBO에서 유지되지 않을 위험 평가 & PASS/YELLOW/HOLD/RED & 강함 \\
Model E. 시장 접근성 모형 & 실제 영입 가능성 평가 & available/conditional/blocked & 제외 조건 \\
Model F. 의료/실행 검증 모형 & 점수 상위 후보의 실행 불가능성 보류 & PASS/YELLOW/HOLD/RED & 제외 조건 \\
\bottomrule
\end{longtable}

\subsection{Team Need Mining Model}
Team Need Mining Model은 선수를 직접 고르는 모델이 아니라, 후보에게 요구되는 조건을 정의하는 모델이다. SSG의 공격 약점은 단순 외야 홈런 부족보다 RHP game-script, run-kill avoidance, extra-out resilience에서 더 구체적으로 나타난다. 투수 약점은 구위 총량보다 traffic 이후 볼넷과 장타 억제, 5이닝 이상을 버티는 선발 안정성에서 나타난다.

\subsection{Historical KBO Success/Failure Model}
이 모델은 과거 KBO 외국인 선수의 입단 전 데이터를 바탕으로 성공/실패 패턴 유사도를 추정한다. 타자 모델은 pre-KBO Savant feature를 사용한 Ridge Logistic Regression 계열로 구성했고, repeated stratified CV에서는 성공 AUC 0.833, 실패 AUC 0.738을 기록했다. 다만 표본이 22행으로 작고 같은 선수가 여러 시즌 반복되므로, 선수 단위 그룹 검증에서는 성공 AUC가 0.650까지 낮아졌다. 따라서 개별 선수 값은 실제 계약 성공확률이 아니라 후보 비교를 위한 모델 지지 점수와 실패 경고 점수로만 해석한다. 투수 모델은 pre-KBO MiLB feature를 기반으로 한 sparse logistic diagnostic model로 구성했으나 AUC 0.603에 그쳐 최종 추천 모델이 아니라 진단 신호로 제한했다.

\begin{table}[H]
\centering
\caption{모델 누수 및 보정 검증 요약}
\label{tab:leakage_calibration_audit}
\small
\begin{tabularx}{\linewidth}{L{0.28\linewidth}L{0.27\linewidth}L{0.35\linewidth}}
\toprule
검증 항목 & 결과 & 해석 \\
\midrule
label/outcome 피처 누수 & 발견 없음 & success, failure, first-KBO 결과, renewal, player id 계열은 학습 피처에서 제외 \\
후보 입력 누수 & 발견 없음 & 후보 입력은 recent Savant/MiLB 구조화 지표로 매핑 \\
train-candidate ID overlap & 2명 발견, 최종 후보와는 무관 & Jared Young, Patrick Wisdom이 후보풀에 남았으나 최종 board에는 없음 \\
선수 단위 그룹 검증 & success AUC 0.650, failure AUC 0.500 & full-fit 확률값을 절대 성공확률처럼 쓰면 과신 위험이 큼 \\
보고서 처리 & 확률 표기 제거 & 모델 지지 점수, 실패 경고 점수, 접촉 순위로만 사용 \\
\bottomrule
\end{tabularx}
\end{table}

\subsection{Similarity/Archetype Matching Model}
Similarity/Archetype Matching은 성적 총량이 아니라 역할 유사성을 본다. 예를 들어 Matos는 단순 파워보다 contact floor와 외야/DH role 안정성에서 SSG 필요 유형과 닮았고, Fleming은 에이스형 구위보다 traffic-command starter 유형에 가깝다.

\subsection{KBO Translation Risk Model}
KBO Translation Risk는 MLB/MiLB 성과가 KBO에서 유지되지 않을 위험을 평가한다. 타자는 K\%, chase, whiff, 변화구 대응, run-kill risk를 본다. 투수는 BB/9, HR/9, starter continuity, pitch-quality 확인 필요성을 본다.

\subsection{시장 접근성과 의료/실행 검증}
시장 접근성 모형은 40-man 상태, DFA/outright, minor contract, salary signal, option status를 평가한다. 의료/실행 검증 모형은 초기 모델 점수가 높아도 즉시 전력 가능성이 낮으면 최종 접촉 후보군에서 보류한다. Carson Spiers와 Brian Van Belle이 그 사례다.

\subsection{앙상블 weight}
타자와 투수는 모델 신뢰도가 다르므로 weight도 다르게 둔다. 타자는 supervised success/failure model의 성능이 비교적 강해 큰 weight를 주고, 투수는 classifier 안정성이 낮아 cross-model consensus와 실행 검증을 더 강하게 반영한다.

\begin{longtable}{L{0.10\linewidth}L{0.34\linewidth}C{0.10\linewidth}L{0.33\linewidth}}
\caption{모델 weight}
\label{tab:model_weights}\\
\toprule
슬롯 & 모듈 & 가중치 & 근거 \\
\midrule
\endfirsthead
\toprule
슬롯 & 모듈 & 가중치 & 근거 \\
\midrule
\endhead
타자 & historical success/failure classifier & 0.40 & 타자 AUC가 높아 주요 base learner로 사용 \\
타자 & SSG fit translation pipeline & 0.25 & SSG 2026 상황별 약점과 KBO 번역 반영 \\
타자 & KBO adaptation filter & 0.15 & contact, chase, 변화구 대응 위험 반영 \\
타자 & market inefficiency feature model & 0.15 & AAA 장점과 MLB 결함의 translation gap 탐색 \\
타자 & cross-model consensus & 0.05 & 여러 base learner 반복 등장 보너스 \\
투수 & historical success/failure classifier & 0.05 & AUC 0.603으로 진단 신호로만 사용 \\
투수 & SSG fit translation pipeline & 0.25 & 선발 이닝, damage control, KBO 번역 반영 \\
투수 & KBO adaptation filter & 0.25 & BB/9, HR/9, starter continuity 반영 \\
투수 & market inefficiency feature model & 0.20 & fringe starter 시장에서 command와 이닝 가치 탐색 \\
투수 & cross-model consensus & 0.25 & 투수 supervised model 약점을 여러 독립 신호로 보완 \\
\bottomrule
\end{longtable}

\subsection{성능 평가}
모델별 성능과 사용 강도는 표 \ref{tab:model_perf}와 같다.

\begin{table}[H]
\centering
\caption{모델 성능 및 사용 강도}
\label{tab:model_perf}
\small
\begin{tabularx}{\linewidth}{L{0.27\linewidth}C{0.12\linewidth}C{0.12\linewidth}L{0.25\linewidth}L{0.13\linewidth}}
\toprule
모델 & 대상 & 표본 수 & 주요 성능 & 사용 강도 \\
\midrule
hitter success classifier & 타자 & 22 & success AUC 0.833 & High \\
hitter failure classifier & 타자 & 22 & failure AUC 0.738 & High \\
pitcher success diagnostic & 투수 & 49 & AUC 0.603 & Low/Diagnostic \\
pitcher warning model & 투수 & 49 & AUC 0.603 & Low/Diagnostic \\
\bottomrule
\end{tabularx}
\end{table}

\section{결과}

\subsection{SSG 보강 문제의 데이터마이닝 결과}
SSG의 타자 보강 문제는 장타 총량이 아니라 우투 선발 경기에서 외야/DH가 이닝을 다시 열지 못하는 구조다. 따라서 후보는 단순 장타형보다 contact floor, two-strike survival, run-kill avoidance를 동시에 충족해야 한다. 투수 보강 문제는 strikeout ceiling보다 traffic 이후 볼넷과 장타를 억제하고 5이닝 이상을 버티는 선발 안정성이다.

\subsection{타자 후보 현실성 반영 Top 3}
타자 data-mining ranking은 과거 KBO 외국인 타자 성공/실패 패턴을 학습한 모델의 지지 신호와 실패 경고 신호를 기준으로 계산했다. 이 값은 calibration된 절대 확률이 아니라 후보 간 우선순위와 위험 방향을 비교하기 위한 모델 산출값이다. 최종 단계에서는 여기에 현실성 검증 단계를 추가했다. 즉, 40-man 또는 최근 cash trade/claim/selected contract로 구단 통제력이 강한 선수, age 25 이하로 MLB 재도전 가치가 큰 선수, OF/DH 역할과 맞지 않는 선수는 모델 신호가 좋아도 최종 접촉 후보군에서 제외했다.

\begin{table}[H]
\centering
\caption{타자 후보 현실성 반영 Top 3}
\label{tab:hitter_raw}
\small
\begin{tabularx}{\linewidth}{C{0.12\linewidth}L{0.20\linewidth}C{0.14\linewidth}C{0.14\linewidth}L{0.13\linewidth}L{0.17\linewidth}}
\toprule
접촉 순위 & 선수 & 모델 지지 & 실패 경고 & 표본 상태 & 주요 지지 모듈 \\
\midrule
1 & Will Brennan & 강함 & 낮음 & 소표본 보정 & Recent DFA/outright, Contact Floor, OF \\
2 & Dominic Fletcher & 강함 & 낮음 & 소표본 보정 & Minor contract, OF, independent signal \\
3 & Dylan Carlson & 중상 & 보통 & 충분 & Historical, Switch-hit, OF Depth \\
\bottomrule
\end{tabularx}
\end{table}

\subsection{타자 후보 실행 보정 Top 3}

\begin{figure}[H]
\centering
\includegraphics[width=0.90\linewidth]{\figpath{v2_hitter_contact_priority.png}}
\caption{타자 후보 최종 접촉 우선순위}
\label{fig:hitter_success}
\end{figure}

\begin{longtable}{C{0.09\linewidth}L{0.18\linewidth}L{0.18\linewidth}L{0.25\linewidth}L{0.20\linewidth}}
\caption{타자 후보 실행 보정 Top 3}
\label{tab:hitter_gate}\\
\toprule
접촉 순위 & 선수 & 최종 판단 & 팀 필요 적합성 & 위험 요약 \\
\midrule
\endfirsthead
\toprule
접촉 순위 & 선수 & 최종 판단 & 팀 필요 적합성 & 위험 요약 \\
\midrule
\endhead
1 & Will Brennan & CONTACT 1ST & DFA/outright 이후 40-man 밖으로 내려온 contact-floor OF & medical history watch는 확인 필요하나, Matos보다 KBO행 현실성이 높음 \\
2 & Dominic Fletcher & CONTACT 2ND & non-40man minor contract OF, 독립 후보 신호와 모델이 교차 & 장타 ceiling은 낮아 중심타선 해결사보다 run-kill avoidance 역할 \\
3 & Dylan Carlson & CONTACT 3RD & 스위치히터 외야수로 platoon 및 OF depth 안정화에 적합 & 최근 소속권과 opt-out/release 조건 확인 필요 \\
\bottomrule
\end{longtable}

\subsection{타자 최종 1인: Will Brennan}
타자 부문에서 초기 발견 신호만 보면 Matos와 Nolan Jones도 강하게 보인다. 그러나 현실성 검증 단계를 앞에 두면 최종 접촉 1순위는 Will Brennan으로 바뀐다. Brennan은 28세 외야수이고, 2026년 6월 10일 DFA, 6월 17일 outright 이후 40-man 밖으로 내려온 고접근성 후보다. 통합 후보 보드에서도 모델 지지 신호가 강하고 실패 경고 신호가 낮게 나타나며, SSG가 원하는 run-kill avoidance와 contact-floor OF 조건에 잘 맞는다.

Matos를 내린 이유는 명확하다. Matos는 데이터마이닝 발견 단계에서 강하지만 24세이고, 이미 현금 트레이드로 한 차례 다른 MLB 구단이 권리를 산 이력이 있다. 이는 단순히 성적이 낮아 방출된 베테랑과 다르다. 선수와 구단 모두 MLB 재도전 옵션을 더 선호할 수 있으므로, KBO 대체 외국인 영입의 현실성 기준에서는 최종 Top 3가 아니라 접촉 보류가 맞다.

\begin{table}[H]
\centering
\caption{초기 발견 리드와 최종 접촉 1순위의 분리}
\label{tab:matos_market_gate}
\small
\begin{tabularx}{\linewidth}{L{0.24\linewidth}L{0.36\linewidth}L{0.30\linewidth}}
\toprule
항목 & 초기 발견 신호 & 현실성 검증 적용 후 \\
\midrule
Matos & 초기 발견 리드, low-K contact floor & age 24, MLB reset value 때문에 HOLD \\
Nolan Jones & power/OBP 발견 신호 & recent cash trade와 salary signal 때문에 HOLD \\
Will Brennan & 모델 점수와 contact floor 모두 강함 & recent DFA/outright로 CONTACT 1ST \\
Dominic Fletcher & 모델 점수와 독립 후보 신호가 교차 & non-40man minor contract로 CONTACT 2ND \\
\bottomrule
\end{tabularx}
\end{table}

\subsection{타자 후보별 차별점}
\begin{table}[H]
\centering
\caption{타자 후보 차별점}
\label{tab:hitter_profiles}
\small
\begin{tabularx}{\linewidth}{L{0.18\linewidth}L{0.26\linewidth}L{0.28\linewidth}L{0.17\linewidth}}
\toprule
선수 & 강점 & 반대 논리 & 최종 판단 \\
\midrule
Will Brennan & recent DFA/outright, contact floor, 낮은 실패 경고 & medical history watch 확인 필요 & CONTACT 1ST \\
Dominic Fletcher & non-40man minor contract, OF 안정성, 독립 후보 신호와 교차 & 장타 ceiling은 높지 않음 & CONTACT 2ND \\
Dylan Carlson & 모델 Top 3를 통과한 스위치히터 안정형 OF & 최근 소속권/release 조건 확인 필요 & CONTACT 3RD \\
Luis Matos & 모델상 contact floor와 강한 발견 신호 & 24세 MLB 재도전 가치가 큼 & 접촉 보류 \\
Nolan Jones & 출루와 장타의 결합, 강한 hard-hit signal & cash trade 및 salary/cost-share 확인 필요 & COST HOLD \\
Jack Suwinski & BB\%와 barrel 기반 upside & 실패 경고와 K\% 위험으로 model warning & WATCHLIST \\
\bottomrule
\end{tabularx}
\end{table}

\subsection{투수 후보 현실성 반영 Top 3}
투수는 supervised classifier 안정성이 낮기 때문에, 초기 모델 점수보다 현재 계약 상태와 선발 지속성을 먼저 본다. Bryse Wilson은 6월 18일 Phillies가 contract를 selected한 이후 현실성 반영 단계에서 hold로 내려간다.

\begin{table}[H]
\centering
\caption{투수 후보 현실성 반영 Top 3}
\label{tab:pitcher_raw}
\small
\begin{tabularx}{\linewidth}{C{0.12\linewidth}L{0.22\linewidth}C{0.14\linewidth}L{0.40\linewidth}}
\toprule
현실 순위 & 선수 & Score & 주요 지지 모듈 \\
\midrule
1 & Josh Fleming & 80.9 & non-40man, AAA starter load, BB/9 1.36, HR/9 0.51 \\
2 & Keegan Thompson & 76.0 & recent DFA/outright, starter load, HR/9 0.56 \\
3 & Kolby Allard & 75.5 & LHP, repeated DFA/outright/minor deal, HR/9 0.36 \\
\bottomrule
\end{tabularx}
\end{table}

\subsection{투수 후보 실행 보정 Top 3}

\begin{figure}[H]
\centering
\includegraphics[width=0.90\linewidth]{\figpath{v2_pitcher_gate_adjustment.png}}
\caption{투수 후보 초기 순위와 실행 보정 순위}
\label{fig:pitcher_gate}
\end{figure}

\begin{longtable}{C{0.09\linewidth}L{0.18\linewidth}L{0.18\linewidth}L{0.25\linewidth}L{0.20\linewidth}}
\caption{투수 후보 실행 보정 Top 3}
\label{tab:pitcher_gate}\\
\toprule
접촉 순위 & 선수 & 최종 판단 & 팀 필요 적합성 & 위험 요약 \\
\midrule
\endfirsthead
\toprule
접촉 순위 & 선수 & 최종 판단 & 팀 필요 적합성 & 위험 요약 \\
\midrule
\endhead
1 & Josh Fleming & CONTACT 1ST & 좌완 선발/스윙맨으로 traffic-command starter 조건에 가장 근접 & 에이스형 구위가 아니라 안정화형 선발로 정의해야 함 \\
2 & Keegan Thompson & CONTACT 2ND & recent DFA/outright high-access, starter load, HR 억제 신호 & K/9 5.29라 swing-and-miss ceiling은 낮음 \\
3 & Kolby Allard & CONTACT 3RD & 좌완, repeated DFA/outright, 낮은 HR/9 & BB/9 3.20과 최근 소속권 확인 필요 \\
\bottomrule
\end{longtable}

\subsection{투수 최종 1인: Josh Fleming}
투수 부문은 타자보다 supervised classifier의 안정성이 낮기 때문에 초기 모델 점수만으로 결론을 내리지 않았다. command, starter floor, 의료/실행 검증, 계약 접근성을 함께 반영했고, 이 과정을 통과한 최종 접촉 1순위는 Josh Fleming이다. Fleming은 좌완 선발/스윙맨 후보이며, 낮은 HR/9, 관리 가능한 BB/9, minor contract 접근성이 결합된다. SSG 내 역할은 traffic-command starter, swingman/spot starter, HR/BB damage controller다.

Fleming을 에이스형 구위 후보로 포장하면 위험하다. 이 후보의 핵심 가치는 압도적 탈삼진이 아니라 traffic 이후 경기 붕괴를 줄이는 안정화 기능이다. 따라서 협상 전에는 최근 선발 빌드업, 구속, 구종 품질, availability를 확인해야 한다.

\subsection{투수 후보별 차별점}
\begin{table}[H]
\centering
\caption{투수 후보 차별점}
\label{tab:pitcher_profiles}
\small
\begin{tabularx}{\linewidth}{L{0.18\linewidth}L{0.26\linewidth}L{0.28\linewidth}L{0.17\linewidth}}
\toprule
선수 & 강점 & 반대 논리 & 최종 판단 \\
\midrule
Josh Fleming & 좌완성, damage control, 시장 접근성 & K/9 upside 제한 & CONTACT 1ST \\
Keegan Thompson & recent DFA/outright, HR 억제, starter load & K/9 5.29로 swing-miss 부족 & CONTACT 2ND \\
Kolby Allard & 좌완, 낮은 HR/9, repeated minor deal & BB/9 3.20과 최근 소속권 확인 필요 & CONTACT 3RD \\
Bruce Zimmermann & K/9 9.98, 67.2 IP/13 GS & HR/9 1.86이 KBO 장타 억제 조건과 충돌 & WATCH \\
Bryse Wilson & 선발 경력과 모델 진단 신호 & 6/18 selected contract 이후 contract blocker & HOLD \\
\bottomrule
\end{tabularx}
\end{table}

\subsection{모델별 지지와 경고}
\begin{table}[H]
\centering
\caption{모델별 후보 지지와 경고}
\label{tab:model_support}
\small
\begin{tabularx}{\linewidth}{L{0.25\linewidth}L{0.25\linewidth}L{0.20\linewidth}L{0.20\linewidth}}
\toprule
모델 & 지지 후보 & 경고 후보 & 이유 \\
\midrule
Team Need Mining & Brennan, Fletcher, Fleming & Suwinski & run-kill avoidance와 traffic-command 조건 충족 여부 \\
Historical Success/Failure & Brennan, Fletcher, Matos & Suwinski & 타자 success/failure ranking score 차이 \\
KBO Translation Risk & Brennan, Carlson, Fleming & Zimmermann, Suwinski & HR/9 및 K\% translation risk \\
시장 접근성 모형 & Brennan, Thompson, Fleming & Matos, Jones, Wilson & age, cash trade, selected contract 등 현실성 차이 \\
의료/실행 검증 모형 & Brennan, Fletcher, Fleming & Spiers, Van Belle & medical hold 여부 \\
\bottomrule
\end{tabularx}
\end{table}

\subsection{검증 후 후보 재분류}
최종 보드 검증에서 가장 중요한 변화는 현실성 검증 단계의 강화다. Matos는 모델상 강하지만 age 24와 MLB 재도전 가치 때문에 최종 접촉 후보군에서 제외했다. Nolan Jones는 모델과 독립 후보 신호가 강하지만 6월 11일 cash trade와 salary signal 때문에 비용/권리 hold로 내렸다. 반대로 Brennan은 6월 17일 outright 이후 40-man 밖으로 내려온 시장 접근성 신호와 모델 지지 신호를 동시에 갖고 있어 최종 접촉 1순위로 올라간다.

\begin{table}[H]
\centering
\caption{검증 후 재분류 결정}
\label{tab:validation_reclass}
\small
\begin{tabularx}{\linewidth}{L{0.18\linewidth}L{0.18\linewidth}L{0.23\linewidth}L{0.30\linewidth}}
\toprule
선수 & 기존 위치 & 검증 후 위치 & 결정 근거 \\
\midrule
Will Brennan & model/market candidate & 최종 타자 접촉 1순위 & 6/17 outright, OF role fit, 낮은 실패 경고 \\
Luis Matos & 발견 단계 lead & 접촉 보류 & age 24와 MLB 재도전 가치 때문에 즉시 KBO행 현실성이 낮음 \\
Nolan Jones & power/OBP 발견 신호 & cost/rights hold & 6/11 cash trade와 salary signal 때문에 대체 외국인 비용 구조 위험 \\
Bryse Wilson & 투수 backup & contract blocker hold & 6/18 selected contract 이후 current 40-man/MLB access blocker \\
\bottomrule
\end{tabularx}
\end{table}

\section{논의}

\subsection{왜 단순 성적 순위와 다른가}
단순 성적 순위는 KBO 번역 위험, 계약 접근성, 의료 상태를 반영하지 않는다. 본 보고서는 data-mining score와 최종 접촉 순위를 분리했기 때문에, 모델 신호는 좋지만 실제 영입 가능성이 낮은 선수를 설명할 수 있다. Matos는 초기 발견 단계에서 강하지만 age 24와 MLB 재도전 가치 때문에 hold로 내려간다. Nolan Jones는 강한 파워/출루 신호가 있지만 cash trade와 salary signal 때문에 비용/권리 hold다. Bryse Wilson은 투수 진단 후보였지만 selected contract 이후 contract blocker로 내려간다.

\subsection{실제 영입 협상에서 확인할 조건}
\begin{table}[H]
\centering
\caption{협상 전 확인 조건}
\label{tab:negotiation_checks}
\small
\begin{tabularx}{\linewidth}{L{0.32\linewidth}L{0.26\linewidth}L{0.32\linewidth}}
\toprule
조건 & 대상 & 의사결정 영향 \\
\midrule
최근 outright/release 권리 & Brennan, Fletcher, Carlson & 즉시 접촉 가능성 및 KBO행 수용 여부 \\
잔여 보장액과 buyout & Jones, Suwinski & 정규 외국인 계약 또는 대체 외국인 비용 구조 충족 여부 \\
young MLB reset value & Matos, Lugo & KBO행보다 MLB 재도전을 택할 위험 \\
최근 구속/구종 품질 & Fleming, Thompson, Allard & KBO starter floor 확인 \\
\bottomrule
\end{tabularx}
\end{table}

\subsection{Plan A / Plan B / Plan C}
\begin{table}[H]
\centering
\caption{접촉 우선순위}
\label{tab:plans}
\small
\begin{tabularx}{\linewidth}{L{0.14\linewidth}L{0.20\linewidth}L{0.20\linewidth}L{0.20\linewidth}L{0.16\linewidth}}
\toprule
슬롯 & Plan A & Plan B & Plan C & 조건 \\
\midrule
타자 & Will Brennan & Dominic Fletcher & Dylan Carlson & Brennan medical/history 확인 후 Fletcher와 병렬 접촉 \\
투수 & Josh Fleming & Keegan Thompson & Kolby Allard & Fleming 선발 빌드업 부족 시 Thompson/Allard 비교 \\
\bottomrule
\end{tabularx}
\end{table}

\section{한계}

\subsection{투수 classifier 안정성이 낮은 이유}
투수 historical classifier의 안정성이 낮은 가장 큰 이유는 표본 수와 label noise다. 현재 투수 학습 표본은 49명 수준이며, KBO 입단 후 투수 성과는 타자보다 주변 환경의 영향을 더 크게 받는다. ERA와 이닝은 투수 본인의 구위뿐 아니라 수비력, 포수 리드, 불펜 승계주자 처리, 구장, 등판 간격, 시즌 중 역할 변경의 영향을 받는다. 따라서 같은 입단 전 BB/9와 HR/9를 가진 투수라도 KBO에서 받은 역할과 팀 환경에 따라 label이 달라질 수 있다.

두 번째 이유는 역할 이질성이다. 외국인 투수 표본 안에는 확정 선발, 스윙맨, 대체 선발, 불펜 전환 후보가 함께 존재한다. 이들은 모두 투수로 묶이지만 성공 기준은 다르다. 선발투수에게 중요한 starter floor와 third-time-through-order risk는 불펜 후보에게는 같은 의미가 아니다. 이 때문에 투수 classifier는 최종 추천 모델이 아니라 후보의 위험 신호를 읽는 진단 모델로 제한한다.

세 번째 이유는 pitch-quality 결측이다. KBO 번역에서 중요한 구속 추세, pitch shape, release consistency, zone command, ABS 환경에서의 called-strike dependency는 공개 집계 성적만으로 완전히 설명되지 않는다. 그래서 투수 최종 판단은 모델 산출값보다 BB/9, HR/9, 최근 선발 이닝, 의료 확인, 계약 접근성, 최근 구속/구종 확인을 함께 보는 방식으로 처리한다.

\subsection{그 외 한계와 보완}
\begin{longtable}{L{0.26\linewidth}L{0.32\linewidth}L{0.30\linewidth}}
\caption{한계와 보완 방향}
\label{tab:limitations}\\
\toprule
한계 & 영향 & 보완 \\
\midrule
\endfirsthead
\toprule
한계 & 영향 & 보완 \\
\midrule
\endhead
투수 KBO 외국인 표본 수 & 연도, 리그, 역할 변화에 민감 & AUC만 보지 않고 진단 신호로 제한 \\
투수 성과 label noise & ERA/IP가 주변 환경 영향을 받음 & BB/9, HR/9, starter floor를 함께 사용 \\
투수 역할 이질성 & 성공 기준이 흔들림 & 역할 적합성 검증을 별도 적용 \\
pitch-quality 결측 & 공개 집계 성적만으로 설명 부족 & 구속/구종 확인을 협상 전 체크리스트로 둠 \\
모델 보정 & 표본 22명 기반 점수를 절대 확률로 과장할 위험 & 점수보다 순위, 방향성, 검증 통과 여부를 함께 제시 \\
salary/contract 공개 정보 & 잔여 부담액과 buyout이 불확실 & 접촉 전 agent/구단 확인 \\
medical public signal & full medical report 부재 & medical hold를 강한 제외 조건으로 운용 \\
2026 시즌 중 데이터 & SSG weakness rule이 갱신될 수 있음 & 주 단위 재빌드 \\
\bottomrule
\end{longtable}

\section{결론}

\subsection{분석 질문별 요약 답변}
\begin{table}[H]
\centering
\caption{분석 질문별 결론}
\label{tab:rq_answers}
\small
\begin{tabularx}{\linewidth}{C{0.10\linewidth}L{0.80\linewidth}}
\toprule
RQ & 요약 답변 \\
\midrule
RQ1 & SSG의 약점은 단순 장타 부족이 아니라 game-state interaction 문제다. \\
RQ2 & 타자는 이닝 전환형 OF/DH, 투수는 traffic-command starter가 필요하다. \\
RQ3 & 후보는 전체 시장, 후보 생성 모듈, 데이터마이닝 1차 검증, 현실성 검증, 최종 Top 3로 축소했다. \\
RQ4 & 각 모델은 팀 필요, historical success, archetype similarity, KBO translation, 시장 접근성, 의료/실행 검증을 분담한다. \\
RQ5 & 타자 Top 3는 Brennan/Fletcher/Carlson, 투수 Top 3는 Fleming/Thompson/Allard다. Matos는 접촉 보류다. \\
RQ6 & 최종 접촉 1순위는 Will Brennan과 Josh Fleming이다. \\
\bottomrule
\end{tabularx}
\end{table}

\subsection{최종 의사결정표}
\begin{table}[H]
\centering
\caption{최종 후보 3인 및 접촉 1순위}
\label{tab:final_decision}
\small
\begin{tabularx}{\linewidth}{L{0.14\linewidth}L{0.18\linewidth}L{0.18\linewidth}L{0.30\linewidth}L{0.12\linewidth}}
\toprule
슬롯 & 초기 발견 후보 & 최종 접촉 1위 & 최종 후보 3명 & 접촉 1순위 \\
\midrule
외국인 타자 & Luis Matos & Will Brennan & Will Brennan, Dominic Fletcher, Dylan Carlson & Will Brennan \\
외국인 선발투수 & Josh Fleming & Josh Fleming & Josh Fleming, Keegan Thompson, Kolby Allard & Josh Fleming \\
\bottomrule
\end{tabularx}
\end{table}

\subsection{핵심 한 줄}
본 보고서의 결론은 Will Brennan과 Josh Fleming이 단순히 가장 좋은 성적의 선수가 아니라, SSG의 약점 구조, KBO 번역 가능성, 시장 접근성, 비용 및 의료 검증을 함께 통과한 현실형 접촉 우선순위라는 점이다. Matos처럼 모델 점수가 높은 후보도 현실성 검증을 통과하지 못하면 최종 접촉 1순위에서 제외한다.

\appendix
\section{부록}

\subsection{최종 후보 보드}
표 \ref{tab:appendix_board}는 검증 후 최종 후보 보드의 축약본이다. 전체 후보 보드 원자료는 별도 CSV 산출물로 함께 보관했다.

\scriptsize
\begin{longtable}{L{0.13\linewidth}L{0.15\linewidth}C{0.07\linewidth}C{0.08\linewidth}L{0.16\linewidth}L{0.27\linewidth}}
\caption{Final candidate board summary}
\label{tab:appendix_board}\\
\toprule
슬롯 & 선수 & 초기 & 접촉 & 최종 판단 & 최종 이유 \\
\midrule
\endfirsthead
\toprule
슬롯 & 선수 & 초기 & 접촉 & 최종 판단 & 최종 이유 \\
\midrule
\endhead
foreign hitter & Will Brennan & 1 & 1 & CONTACT 1ST & 6/17 outright, high market access, contact-floor OF \\
foreign hitter & Dominic Fletcher & 2 & 2 & CONTACT 2ND & non-40man minor contract, 독립 후보 신호와 모델 교차 \\
foreign hitter & Dylan Carlson & 3 & 3 & CONTACT 3RD & switch-hitting OF, 모델 Top 3 통과 \\
foreign hitter & Luis Matos & D & HOLD & 발견 단계 보류 & 초기 발견 리드였지만 age 24와 MLB reset value 때문에 hold \\
foreign hitter & Nolan Jones & 2 & HOLD & COST HOLD & cash trade와 salary signal 때문에 비용/권리 확인 필요 \\
foreign pitcher & Josh Fleming & 1 & 1 & CONTACT 1ST & 53.0 IP/10 GS, BB/9 1.36, HR/9 0.51 \\
foreign pitcher & Keegan Thompson & 2 & 2 & CONTACT 2ND & recent DFA/outright, HR/9 0.56 \\
foreign pitcher & Kolby Allard & 3 & 3 & CONTACT 3RD & 좌완, repeated minor deal, HR/9 0.36 \\
foreign pitcher & Bryse Wilson & 15 & HOLD & CONTRACT HOLD & 6/18 selected contract 이후 active MLB access blocker \\
\bottomrule
\end{longtable}

\clearpage
\begin{thebibliography}{9}
\bibitem{crispdm}
Chapman, P., Clinton, J., Kerber, R., Khabaza, T., Reinartz, T., Shearer, C., and Wirth, R. (2000). CRISP-DM 1.0: Step-by-step data mining guide.

\bibitem{mlb_savant}
MLB Advanced Media. Baseball Savant Statcast Search and Player Pages.

\bibitem{mlb_transactions}
MLB Advanced Media. MLB.com player transaction logs and team transaction pages, 2025-2026.

\bibitem{mlbtr}
MLB Trade Rumors and ESPN transaction pages. Minor-league contract, outright, and selected-contract reports, 2025-2026.

\bibitem{statiz}
STATIZ. KBO team, player, and situational statistics.

\bibitem{kbo_foreign}
Korea Baseball Organization. KBO league registration and foreign player regulations.
\end{thebibliography}

\end{document}
